\documentclass[12pt,aspectratio=169]{beamer}
\input{header_beamer}
\usetheme{Boadilla}
\usecolortheme{beaver}
\setbeamertemplate{page number in head/foot}{}
\setbeamertemplate{navigation symbols}{}
\title{Lecture-25: Poisson processes: Conditional distribution}
\date{}

\begin{document}
\pagestyle{empty}
\maketitle


\begin{frame}{Joint conditional distribution of points in a finite window}
Let $\sX =\R^d$ be a $d$-dimensional Euclidean space, 
and $S:\Omega\to\sX^\N$ be a Poisson point process with intensity measure $\Lambda: \sB(\sX) \to \R_+$ and associated counting process $N:\Omega\to\Z_+^{\sB(\sX)}$.
\begin{prop}<2->
Let $k \in \N$ be any positive integer. 
Consider a Poisson point process $S: \Omega \to \sX^\N$ with intensity measure $\Lambda:\sB(\sX)\to\R_+$,  
a finite partition $A\in \sB(\sX)^k$ that partitions a bounded set $B \in \sB(\sX)$, 
and a vector $n \in \Z_+^k$ that partitions a non negative integer $m \in \Z_+$. 
Then, 
\EQN{
\label{eqn:multi}
P(\set{N(A_1) = n_1, \dots, N(A_k) = n_k}\given\set{N(B) = m}) 
= 
\binom{m}{n_1, \dots, n_k}\prod_{i=1}^k\left(\frac{\Lambda(A_i)}{\Lambda(B)}\right)^{n_i}.
}
\end{prop}
\end{frame}

\begin{frame}
\begin{block}{Proof} 
From the definition of conditional probability and the fact that $\cap_{i=1}^k\set{N(A_i) = n_i} \subseteq \set{N(B) =m}$, we can write the conditional probability on LHS as the ratio 
\begin{equation*}
\frac{P\set{N(A_1) = n_1, \dots, N(A_k) = n_k, N(B) = m}}{P\set{N(B) = m}} = 
\frac{P\set{N(A_1) = n_1, \dots, N(A_k) = n_k}}{P\set{N(B) = m}}.
\end{equation*}
\end{block}
\end{frame}

\begin{frame}
\begin{Proof}[Proof continued]
From the complete independence property and Poisson marginals for the joint distribution of $(N(A_1), \dots, N(A_k))$ for the partition $A \in \sB(\sX)^k$, 
and the fact that the intensity measure sums over disjoint sets, i.e. $\Lambda(B) = \sum_{i=1}^k\Lambda(A_i)$, we can rewrite the RHS of the above equation as 
\begin{align*}
\frac{P\set{N(A_1) = n_1, \dots, N(A_k) = n_k}}{P\set{N(B) = m}}
&= \frac{\prod_{i=1}^ke^{-\Lambda(A_i)}\frac{\Lambda(A_i)^{n_i}}{n_i!}}{e^{-\Lambda(B)}\frac{\Lambda(B)^m}{m!}}\\
&= \binom{m}{n_1, \dots, n_k}\prod_{i=1}^k\left(\frac{\Lambda(A_i)}{\Lambda(B)}\right)^{n_i}.
\end{align*}
\end{Proof}
\end{frame}

\begin{frame}
\begin{block}{Reminder}
Consider a Poisson point process $S:\Omega\to\sX^\N$ with intensity measure $\Lambda: \sB(\sX) \to \R_+$ and counting process $N:\Omega\to\Z_+^{\sB(\sX)}$. 
Let $A \in \sB(\sX)^k$ be a partition for bounded set $B \in \sB(\sX)$.  
\begin{enumerate}[i\_]
\item Defining $p_i \triangleq \frac{\Lambda(A_i)}{\Lambda(B)}$, we see that $(p_1, \dots, p_k) \in \cM([k])$ is a probability distribution. % over partitions $(A_1, \dots, A_k)$. 
We also observe that 
\EQ{
p_i 
= P(\set{N(A_i)=1}\given\set{N(B) = 1}) 
= P(\set{\abs{S \cap A_i} = 1}\given\set{\abs{S \cap B} = 1}).
}
When $N(B)=1$, we can call the point of $S$ in $B$ as $S_1$ without any loss of generality. 
That is, if we call $\set{S_1} = S \cap B$, then we have 
\EQ{
p_i = P(\set{S_1 \in A_i}\given\set{S_1 \in B}).
}
\end{enumerate}
\end{block}
\end{frame}

\begin{frame}
\begin{block}{Reminder}
\begin{enumerate}[i\_]
\item (cntd.) Similarly, when $N(B) = n_i$, 
we call the points of $S$ in $B$ as $S_1, \dots, S_{n_i}$ and denote $S\cap B = \set{S_1, \dots, S_{n_i}}$. 
For this case,  we observe
\eq{
&P(\set{N(A_i)=n_i}\given\set{N(B) = n_i}) 
= P(\set{\abs{S \cap A_i} = n_i}\given\set{\abs{S \cap B} = n_i}) = p_i^{n_i}\\
&= P(\cap_{j=1}^{n_i}\set{S_j \in A_i}\given \set{\set{S_1, \dots, S_{n_i}} = S\cap B}) 
= \prod_{j=1}^{n_i}P(\set{S_j \in A_i}\given \set{S_j \in B}). 
}

\item We can rewrite the Equation~\eqref{eqn:multi} as a multinomial distribution, where 
\EQ{
 P(\set{N(A_1) = n_1, \dots, N(A_k) = n_k}\given\set{N(B) = m}) = {m \choose n_1, \dots, n_k} p_1^{n_1} \dots p_k^{n_k}.
} 
\end{enumerate}
\end{block}
\end{frame}

\begin{frame}
\begin{block}{Reminder}
\begin{enumerate}[i\_]
\setcounter{enumi}{2}
\item For any finite set $F\subseteq \N$ of size $m \in \Z_+$ and $n\in\Z_+^k$ a partition of $m$, 
we define $\cP_k(F, n)$ to be the collection of all $k$-partitions $E \in \cP(\N)^k$ of $F$ such that $\abs{E_i} = n_i$ for $i \in [k]$. 
Then, the multinomial coefficient accounts for number of partitions of $m$ points into sets with $n_1,  \dots, n_k$ points. 
That is, 
\EQ{ 
{m \choose n_1, \dots, n_k} = \abs{\cP_k([m], n)}.
}

\item %Defining $E_i \triangleq S \cap A_i$ and $E = S \cap A$, 
Recall that the event $\set{N(A_i) = n_i} = \set{\abs{S \cap A_i} = n_i}$. 
Hence, we can write 
\eq{
&P(\cap_{i=1}^k\set{\abs{S \cap A_i} = n_i}\given \set{\abs{S \cap B} = m}) 
= {m \choose n_1, \dots, n_k} p_1^{n_1} \dots p_k^{n_k}\\
&= \sum_{E \in \cP_k(S \cap B, n)}\prod_{i=1}^k\prod_{S_j\in E_i}P(\set{S_j \in A_i}\given \set{S_j \in B}).
}
\end{enumerate}
\end{block}
\end{frame}

\begin{frame}
\begin{block}{Reminder}
\begin{enumerate}[i\_]
\setcounter{enumi}{4}
\item 
When $N(B) = m$, we denote $S\cap B$ by $F = \set{S_1, \dots, S_m}$ without any loss of generality. 
We further observe that when $N(A_i) = n_i$ for all $i \in [k]$, 
then $(S \cap A_1, \dots, S \cap A_k) \in \cP_k(S\cap B,n)$. 
Therefore, we can re-write the event 
\EQ{
\cap_{i=1}^k\set{N(A_i) = n_i} 
= \cap_{i=1}^k\set{\abs{S \cap A_i} = n_i} 
= \cup_{E \in \cP_k(S\cap B, n)}(\cap_{i=1}^k\set{S \cap A_i = E_i}).
} 
That is, we can write the conditional probability conditioned on $S\cap B = F$, as 
\eq{
&P(\cap_{i=1}^k\set{N(A_i) = n_i} \given \set{N(B) = m}) \\ 
&= \sum_{E \in \cP_k(F,n)}P(\cap_{i=1}^k\set{S \cap A_i = E_i}\given\set{S\cap B = F})\\
&= \sum_{E \in \cP_k(F, n_1, \dots, n_k)}P(\cap_{i=1}^k\cap_{S_j \in E_i}\set{S_j \in A_i}\given\set{S\cap B = F}).
}
\end{enumerate}
\end{block}
\end{frame}

\begin{frame}
\begin{block}{Reminder}
\begin{enumerate}[i\_]
\setcounter{enumi}{5}
\item   
Equating the RHS of the above equation term-wise, we obtain that conditioned on each of these points falling inside the window $B$, 
the conditional probability of each point falling in partition $A_i$ is independent of all other points and given by $p_i$. 
That is, we have 
\eq{
 P(\cap_{i=1}^k\cap_{S_j \in E_i}\set{S_j \in A_i}\given\set{S\cap B = F}) 
 &= \prod_{i=1}^k\prod_{S_j\in E_i}P(\set{S_j \in A_i}\given \set{S_j \in B}) \\ 
 &= \prod_{i=1}^kp_i^{n_i} =  \prod_{i=1}^k\left(\frac{\Lambda(A_i)}{\Lambda(B)}\right)^{n_i} .
} 
It means that given $m$ points in the window $B$, the location of these points are independently and identically distributed in $B$ according to the distribution $\frac{\Lambda(\cdot)}{\Lambda(B)}$. 
\end{enumerate}
\end{block}
\end{frame}

\begin{frame}
\begin{block}{Reminder}
\begin{enumerate}[i\_]
\setcounter{enumi}{6}
\item If the Poisson process is homogeneous, the distribution is uniform over the window $B$. 
\item For a Poisson process with intensity measure $\Lambda$ and any bounded set $A \in \sB$, 
the number of points $N(A)$ in the set $A$ is a Poisson random variable with parameter $\Lambda(A)$. 
Given the number of points $N(A)$, the location of all the points in $S\cap A$ are \iid with density $\frac{\lambda(x)}{\Lambda(A)}$ for all $x \in A$. 
\end{enumerate}
\end{block}
\end{frame}

\begin{frame}
\begin{block}{Reminder: Simulating a homogeneous Poisson point process} 
If we are interested in simulating a two dimensional homogeneous Poisson point process with density $\lambda$ in a uniform square $A = [0,1]\times[0,1]$. 
Then, we first generate the random variable $N(A):\Omega\to\Z_+$ that takes value $n$ with probability $e^{-\lambda}\frac{\lambda^n}{n!}$. 
Next, for each of the $N(A)=n$ points, we generate the location $(X_i, Y_i) \in \R^2$ uniformly at random. 
That is, $X:\Omega\to[0,1]^n$ and $Y:\Omega\to[0,1]^n$ are independent \iid uniform sequences.
\end{block}
\end{frame}

\begin{frame}
\begin{cor} 
For a homogeneous Poisson point process on the half-line with ordered set of points $\tilde{S} = (S_{(n)} \in \R_+: n \in \N)$, 
we can write the conditional density of ordered points $(S_{(1)}, \dots, S_{(k)})$ given the event $\set{N_t = k}$ as the ordered statistics of \iid uniformly distributed random variables. 
Specifically, we have 
\EQ{
f_{S_{(1)}, \dots, S_{(k)}\given N_t = k}(t_1, \dots, t_k) = k!\frac{\SetIn{0 < t_1 \le \dots \le t_k \le t}}{t^k}. 
}
\end{cor}
\end{frame}

\begin{frame}
\begin{block}{Proof}
Given $\set{N_t = k}$, we can denote the points of the Poisson process in $(0,t]$ by $S_1, \dots, S_k$. 
From the above remark, we know that $S_1, \dots, S_k$ are \iid uniform in $(0, t]$, 
conditioned on the number of points $N_t = k$. 
Hence, we can write 
\eq{
F_{S_1, \dots, S_k\given N_t = k}(t_1, \dots, t_k) 
&= P(\cap_{i=1}^k\set{S_i \in (0, t_i]}\given \set{N_t=k}) \\
&= \prod_{i=1}^kP(\set{S_i \in (0, t_i]}\given \set{S_i \in (0, t]}) 
= \prod_{i=1}^k\frac{t_i}{t}\SetIn{0 < t_i \le t}.
}  
Therefore, for $0 < t_1 < \dots < t_k < 1$ and $(h_1, \dots, h_k)$ sufficiently small, we have 
\EQ{
P\Big(\cap_{i=1}^k\set{S_i \in (t_i, t_i+h_i]}\Big) = \prod_{i=1}^k\frac{h_i}{t}.
}
\end{block}
\end{frame}

\begin{frame}
\begin{Proof}[Proof continued]
Since $(S_1, \dots, S_k)$ are conditionally independent given $S\cap A = \set{S_1, \dots, S_k}$, 
it follows that any permutation $\sigma: [k] \to [k]$, the conditional joint distribution of $(S_{\sigma(1)}, \dots, S_{\sigma(k)})$ is identical to that of $(S_1, \dots, S_k)$ given $S\cap A = \set{S_1, \dots, S_k}$. 
Further, we observe that the order statistics of $(S_{\sigma(1)}, \dots, S_{\sigma(k)})$ is identical to that of $(S_{1}, \dots, S_{k})$.  
Therefore, we can write the following equality for the events
\EQ{
\cap_{i=1}^k\set{S_{(i)}\in (t_i, t_i+h_i]} 
= \cup_{\sigma: [k] \to [k] \text{ permutation}}\cap_{i=1}^k\set{S_{\sigma(i)}\in (t_i, t_i+h_i]}. 
}
The result follows since the number of permutations $\sigma: [k] \to [k]$ is $k!$. 
 %P(\cap_{i=1}^k\set{S_{(i)} \in (0, t_i]}\given \set{N_t=k})  
\end{Proof}
\end{frame}





\end{document}
